\textbf{1.} Pour tous réels $a$, $b$, $c$,
\[
2(a^2 + b^2 + c^2) - 2(ab + bc + ca) = (a-b)^2 + (b-c)^2 + (c-a)^2 \geqslant 0,
\]
donc $ab + bc + ca \leqslant a^2 + b^2 + c^2 = 1$. L'égalité a lieu si et
seulement si $a = b = c$ ; avec $3a^2 = 1$, cela donne exactement les deux
triplets $a = b = c = \frac{1}{\sqrt3}$ et $a = b = c = -\frac{1}{\sqrt3}$.

\textbf{2.} On a $(a+b+c)^2 = a^2 + b^2 + c^2 + 2(ab+bc+ca) = 1 + 2(ab+bc+ca)$, donc
\[
ab + bc + ca = \frac{(a+b+c)^2 - 1}{2} \geqslant \frac{0 - 1}{2} = -\frac12 ,
\]
car un carré est positif.

\textbf{3.} L'égalité a lieu si et seulement si $a + b + c = 0$ (avec
$a^2 + b^2 + c^2 = 1$). Par exemple $(a, b, c) = \left(\frac{1}{\sqrt2}, -\frac{1}{\sqrt2}, 0\right)$ :
on a bien $a^2 + b^2 + c^2 = \frac12 + \frac12 = 1$ et
$ab + bc + ca = -\frac12$. \hfill $\blacksquare$
